        %auto-ignore

\section{Nonlinear Depth-$\mu$P Limit}

When the nonlinearity $\phi$ is nontrivial, the distribution of the final representations $x^L_s$ may be highly non-Gaussian because of complex compositions of $\phi$ and $\phi'$, as one would suspect from known examples of large width limits.
This is indeed the case for finite depth $L$. But in fact, when $L \to \infty$, $\{x^L_s\}_{s\ge 0}$ becomes a Gaussian process again!

The kernel of the limiting GP can be computed in a similar way as in the linear case:
\begin{definition}\label{defn:nonlinearGammaC}
    Define $\Gamma$ and $C$ recursively by
    \begin{align*}
        \Gamma_{t,r,0,b}(p, q)&=\ind_{[(t=r) \wedge (b=0) \wedge (p\geq q)]} \\ 
        &+ \int_0^p \sum_{s=0}^{t-1} V_{\phi'}[C]_{t,s,0}(p')\Gamma_{s,r,1,b}(p',q)\cdot (\Gamma_{t,s,0,1}(p', p') - C_{t,s,0}(p'))\dd p'; \\
           \Gamma_{t,r,1,b}(p, q) &= \ind_{[(t=r) \wedge (b=1) \wedge (p\leq q)]} \\
           &+ \int_p^1 \sum_{s=0}^{t-1} \Gamma_{s,r,0,b}(p',q) \cdot (V_{\phi_c|\phi'}[C]_{t,s,0}\Gamma_{t,s,1,0}(p',p') - V_{\phi'}[C]_{t,s,0}(p') C_{t,s,1}(p'))\dd p';\\
           C_{t,s,a}(p)&=
        %
           \sum_{b\in\{0,1\}}\int_0^1 \Gamma_{t,t',a,b}(p,q) V_{\phi_c} [C]_{t',s',b}(q)\Gamma_{s,s',a,b}(p,q) \dd q.
    \end{align*}
\end{definition}

Here
\begin{align*}
    V_{\phi_c}[C]_{t,s,b}(p) &\defeq \EV \MS(\phi(z)) \MS(\phi(y)) \\ 
    V_{\phi'}[C]_{t,s,b}(p) &\defeq \EV \phi'(z) \phi'(y) \\
    V_{\phi_c | \phi'}[C]_{t,s,b}(p) &\defeq \EV \MS(\phi'(z)) \phi'(y) \\
\end{align*}
where $(z, y) \sim \Gaus(0, C^b(p)|_{\{t,s\}})$.

Then 
\begin{claim}\label{claim:nonlineardepthmuplimit}
    For sufficiently smooth nonlinearity $\phi$, in the $L \to \infty$ limit, the kets $\ket{x^L_s}, s = 0, 1, \ldots,$ converge in distribution as a zero-mean Gaussian process with kernel
    \[\braket{x^L_s}{x^L_t} = C_{t,s,1}(1).\]
    as defined in \Cref{defn:nonlinearGammaC}.
    Thus, for each fixed neuron index $\alpha$, the collection $\{x^L_{\alpha s}\}_{s\ge 0}$ converges in distribution to a zero-mean Gaussian process with kernel $C_{t,s,1}(1)$ in the $n\to\infty$ then $L\to\infty$ limit.
\end{claim}
We frame this as a claim because we do not want to get into the details of what ``sufficiently smooth nonlinearity'' means here, nor do we give a proof. Instead, we give an intuitive justification.

\begin{proof}[Heuristic justification of \Cref{claim:nonlineardepthmuplimit}]
First, in Depth-$\mu$P, we can Taylor expand each block
$$\ket{\phi(W^l_t x^{l-1}_t)} = \phi(\hatket{W^l_0 x^{l-1}_t}) + \phi'(\hatket{W^l_0 x^{l-1}_t}) [\ket{\bar\Delta W^l_t x^{l-1}_t} + \dotket{W^l_0 x^{l-1}_t}] + \cO(L^{-1}).$$
The remainder term thus contributes $\cO(L^{-3/2})$ after accounting for the $L^{-1/2}$ block multiplier. Summing over depth $l \in [0, L]$, by Gronwall's lemma, the remainders of all layers sum up to $\cO(L^{-1/2})$, so we can ignore them.
From here on, we study the linearized blocks $x \mapsto \MS(\phi(W^l_0 x) + \phi'(W^l_0 x) \odot \bar\Delta W^l_t x)$.

Now the key observation is that each of $\ket{x^l_t}$ or $\ket{\delta x^l_t}$ is always equal, up to $\cO(L^{-1/2})$ factor, to a linear combination of $\{\MS \phi(\hatket{W^l_0 x^{l-1}_s}), \hatket{W^{l \trsp}_0 \delta h^l_s}\}_{l, s}$, where each coefficient in this linear combination is $\cO(L^{-1/2})$ and deterministic.
We call such a linear combination a \emph{good linear combination}.
One can see this by an induction argument on $t$.
Indeed, at initialization $t=0$, this claim is trivially true.
Supposing this claim is true for $t$, then it is trivial to see that it remains true at $t+1$ for the backward pass kets $\ket{\delta x^l_t}$.
The only nontrivial part is to show the forward pass for $t+1$.
By the Taylor expansion above, we just need to show that $\phi'(\hatket{W^l_0 x^{l-1}_t}) [\ket{\bar\Delta W^l_t x^{l-1}_t} + \dotket{W^l_0 x^{l-1}_t}]$ is ``well-approximated'' by a good linear combination.
By induction, both $\ket{\bar\Delta W^l_t x^{l-1}_t}$ and $\dotket{W^l_0 x^{l-1}_t}$ are of products between a good linear combination and a term of form $ \phi'(\hatket{W^m_0 x^{m-1}_s})$ for some $s \le t$. Thus $\phi'(\hatket{W^l_0 x^{l-1}_t}) [\ket{\bar\Delta W^l_t x^{l-1}_t} + \dotket{W^l_0 x^{l-1}_t}]$ is of the form $\phi'(\hatket{W^l_0 x^{l-1}_t})\phi'(\hatket{W^m_0 x^{m-1}_s})$ times a good linear combination.
But $\phi'(\hatket{W^l_0 x^{l-1}_t})\phi'(\hatket{W^m_0 x^{m-1}_s})$ only correlates with a single component in this linear combination and is independent with all other components.
Therefore, when summing over depth, the noncorrelated components will experience law of large numbers and one can replace $\phi'(\hatket{W^l_0 x^{l-1}_t})\phi'(\hatket{W^m_0 x^{m-1}_s})$ with its expectation; the correlated components is only a sum of $\cO(L)$ elements each of size $\cO(L^{-3/2})$, so in total they are $o(1)$.
This completes the induction.

In this reasoning, the coefficients of a good linear combination correspond to $\Gamma$ in \Cref{defn:nonlinearGammaC}, and filling in the details of the induction yields the recursive formula in \Cref{defn:nonlinearGammaC}.

Finally, because a good linear combination is a large sum of independent terms, the central limit theorem tells us that $\{\ket{x^l_t}, \ket{\delta x^l_t}\}_{l, t}$ converge in distribution to a Gaussian process.
\end{proof}

\Cref{claim:nonlineardepthmuplimit} is good news for theorists, that we have such a simple form for a fundamental architecture.
At the same time, one may worry that this Gaussian form lacks expressivity.
But in fact, some common architecture or algorithm choices would make the limit non-Gaussian.
For example, the use of adaptive optimizers like Adam or SignSGD.
Or the addition of a nonlinearity before the matrix multiply, i.e., ``prenonlin'', (in addition to the one afterward, ``postnonlin'').

In general, one obtains a stochastic differential equation with McKean-Vlasov elements describing the evolution of $x^l_t$ over depth and time. However, the stochastic integral involved is \emph{not} the usual Ito or Stratonovich integral because the depthwise evolution requires integrating non-adapted process against Brownian motion. Instead, we need to use \emph{Skorohod} integral and the SDE is only defined in the sense of Malliavin calculus.
This is not just a mathematical nitpick; rather, assuming Ito calculus (which amounts to assuming incorrect independence between certain quantities) will lead to the wrong predictions and calculations.
Malliavin calculus is intimately connected to Tensor Programs and we shall develop their relationship as well as the theory of the general infinite-depth limit in a future work.

